\documentclass{article}
\usepackage[utf8]{inputenc}

\title{Design Specification: Knowledge Expiration Tracker}
\author{Daniel Yang, Chad Wong, Raymond Diep, William Lopez, Liam Peak}
\date{\today}

\begin{document}

\maketitle
\newpage

\section{System Overview}
The Knowledge Expiration Tracker is built using a containerized architecture to ensure consistency across development environments. The project utilizes a frontend-backend-database structure managed via Docker Compose.

\section{Frontend (Client-Side)}
The frontend is developed using React.js to provide a responsive user interface for managing study schedules and flashcards.
\begin{itemize}
    \item \textbf{Framework:} React.js
    \item \textbf{UI Library:} Material UI (for responsive design and components)
    \item \textbf{Core Dependencies:} react, react-dom, react-scripts
\end{itemize}

\section{Backend (Server-Side)}
The backend handles the core logic for calculating knowledge expiration dates and managing user accounts.
\begin{itemize}
    \item \textbf{Runtime:} Node.js
    \item \textbf{Framework:} Express.js
    \item \textbf{Core Dependencies:} express, body-parser, cors, dotenv
\end{itemize}

\section{Database}
User data and knowledge items are stored in a relational database to maintain study history and scheduled intervals.
\begin{itemize}
    \item \textbf{System:} PostgreSQL
    \item \textbf{Library:} pg (node-postgres)
\end{itemize}

\section{Security and Authentication}
Secure access to personal study data is managed through token-based authentication.
\begin{itemize}
    \item \textbf{Tool:} jsonwebtoken (JWT)
\end{itemize}

\section{Infrastructure and Containerization}
The entire stack is orchestrated using Docker to ensure all dependencies are correctly loaded in the production image.
\begin{itemize}
    \item \textbf{Tool:} Docker Compose
    \item \textbf{Package Manager:} npm
\end{itemize}

\end{document}
